\chapter{Multiagent Reinforcement Learning for VAMPs}\label{chap:marl}

% \subsection{Executive Summary}

% We develop a MARL framework for decentralized HTAP (formalized as a VAMP) that targets \textbf{empirical low-regret behavior} rather than exact equilibria, because (i) Nash computation is PPAD-hard even in one-shot subcases, and (ii) optimal decentralized partially observable control is NEXP-hard even in cooperative special cases. Our approach follows the structure already drafted: (a) interpret learning as equilibrium approximation in a \emph{policy-population meta-game}, (b) use CTDE actor--critic training (MAPPO/PPO-style) plus sequence models (causal transformers) for variable-length theorem-proving observations, and (c) use population-based training/PSRO-style iteration to reduce overfitting and approximate meta-game CCE. Rigorous guarantees are necessarily weak: we can prove CCE convergence in repeated normal-form meta-games under no-regret assumptions, while in the underlying VAMP these guarantees degrade to ``restricted-class'' statements depending on function approximation, exploration, and oracle quality. Claims requiring heavy additional machinery (e.g., finite-sample deep PPO guarantees in POSGs) are marked \textbf{TODO}.

% \subsection{Formal MARL Problem Statement for VAMPs}

% We assume a VAMP \(\mathsf{V}\) (finite-horizon \(H\) or discounted \(\gamma\in(0,1)\)) with agent set \(\mathcal N\), world state \(w_t\in\mathcal W\), per-agent observation \(o_t^\alpha\in\mathcal O^\alpha\), and joint action \(a_t=(a_t^\alpha)_{\alpha\in\mathcal N}\), where each \(a_t^\alpha\) factors into compute actions (prove/conjecture/publish) and market actions (trade/offer/cancel) subject to admissibility/collateral constraints.

% \subsubsection{Policies, histories, and policy classes}

% Let \(h_t^\alpha=(o_0^\alpha,a_0^\alpha,\dots,o_t^\alpha)\in(\mathcal O^\alpha\times \mathcal A^\alpha)^*\) denote private action--observation history.

% A stochastic policy is
% \[
% \pi^\alpha_\theta(\cdot \mid h_t^\alpha)\in \Delta(\mathcal A^\alpha),
% \]
% parameterized by \(\theta\in\Theta^\alpha\) (e.g., neural parameters). We write the joint policy \(\pi_\theta = (\pi^\alpha_{\theta^\alpha})_{\alpha\in\mathcal N}\).

% Because the full pure-strategy space is exponential in \(H\), our learning guarantees will always be relative to a restricted class \(\Pi^\alpha=\{\pi_{\theta^\alpha}^\alpha:\theta^\alpha\in\Theta^\alpha\}\).

% \subsubsection{CTDE (centralized training with decentralized execution)}

% In CTDE we train decentralized actors \(\pi^\alpha_{\theta^\alpha}(a^\alpha\mid h^\alpha)\) that condition only on local information at execution time, while permitting a critic to access richer state during training (e.g., joint public state, market summaries, cross-agent statistics). This is the standard CTDE paradigm in cooperative MARL and underlies algorithms such as MADDPG and MAPPO.

% We formalize CTDE by specifying a centralized value function class
% \[
% V_\psi : \mathcal W \to \mathbb R
% \quad\text{or}\quad
% Q_\psi : \mathcal W\times\mathcal A \to \mathbb R,
% \]
% used only in the gradient estimator / advantage computation.

% \subsubsection{Objectives and evaluation metrics}

% \paragraph{Per-agent return.}
% For each agent \(\alpha\), define discounted return
% \[
% J^\alpha(\pi) \;:=\; \mathbb E_\pi\Big[\sum_{t=0}^{H-1}\gamma^t\, r^\alpha(w_t,a_t)\Big],
% \]
% with the understanding that \(r^\alpha\) may include both task-related utility and market profit/loss.

% \paragraph{Equilibrium-approximation target (meta-game view).}
% In general-sum VAMPs, we do not aim for Nash of the full POSG. Instead we aim for \textbf{low regret} in a repeated interaction among a policy population (formalized below). This is consistent with the chapter draft framing.

% \paragraph{Regret.}
% Over \(T\) training episodes (or ``meta-rounds''), agent \(\alpha\) suffers external regret w.r.t.\ policy class \(\Pi^\alpha\):
% \[
% \mathrm{Reg}^\alpha(T)
% := \max_{\pi'\in\Pi^\alpha}\ \sum_{k=1}^T J^\alpha(\pi',\pi^{-\,\alpha}_k)\ -\ \sum_{k=1}^T J^\alpha(\pi^\alpha_k,\pi^{-\,\alpha}_k),
% \]
% where \(\pi_k\) is the profile used in episode \(k\). Low regret means \(\mathrm{Reg}^\alpha(T)=o(T)\).

% \paragraph{Exploitability.}
% For two-player zero-sum special cases, define exploitability of a profile \(\pi\) as the sum of best-response gaps. (We use this mainly for diagnostics in structured settings; see PSRO below.)

% \paragraph{Sample complexity notions (conceptual).}
% Let \(\widehat\pi_T\) denote the learned output after \(T\) environment transitions. A canonical PAC-style notion is:
% \[
% N(\varepsilon,\delta) := \min\Big\{T:\ \Pr\big(J^\alpha(\widehat\pi_T)\ge J^\alpha(\pi^\star)-\varepsilon,\ \forall \alpha\big)\ge 1-\delta\Big\},
% \]
% for a reference comparator \(\pi^\star\) (e.g., best in class, or best response to a fixed mixture). In deep MARL for POSGs, finite-sample bounds of this kind generally require strong realizability/mixing assumptions and are rarely tight; we therefore use sample complexity primarily as an experimental reporting metric (environment steps to reach a given regret/utility threshold). \textbf{TODO:} cite a primary finite-sample theory result tailored to deep on-policy actor--critic in partial observability if you want formal non-asymptotic bounds.

% \subsection{Algorithmic Frameworks for VAMP-MARL}

% \subsubsection{On-policy actor--critic with PPO/MAPPO}

% We use PPO-style clipped policy-gradient updates as the core optimizer, extended to multiagent CTDE as MAPPO-style training (shared or separate parameters, centralized critic, decentralized actors). Empirically, PPO-style multiagent training is competitive and widely used.

% For each agent \(\alpha\), PPO optimizes a surrogate objective
% \[
% \max_{\theta^\alpha}\ \mathbb E\Big[\min\big(r_t(\theta^\alpha)\,\widehat A_t,\ \mathrm{clip}(r_t(\theta^\alpha),1-\epsilon,1+\epsilon)\,\widehat A_t\big)\Big],
% \]
% where \(r_t(\theta^\alpha)=\frac{\pi_{\theta^\alpha}^\alpha(a_t^\alpha\mid h_t^\alpha)}{\pi_{\theta^\alpha_{\mathrm{old}}}^\alpha(a_t^\alpha\mid h_t^\alpha)}\) and \(\widehat A_t\) is an advantage estimate from the centralized critic.

% \paragraph{Masking and admissibility.}
% In VAMPs, market actions are constrained by admissibility/collateral predicates; we enforce this via action masking in the actor's output distribution. This is an implementation-level constraint that preserves correctness of sampling but complicates theoretical analysis (masked policy gradients). \textbf{TODO:} cite/derive a masked-policy PPO convergence statement.

% \subsubsection{Population-based training and PSRO-style population dynamics}

% A single self-play equilibrium target is brittle in general-sum MARL. We therefore maintain a population of policies, sampling opponents from the population (``league training''). This matches the draft's ``population-based training approximates PSRO.''

% Policy-Space Response Oracles (PSRO) provides the game-theoretic template: iteratively compute approximate best responses to mixtures of opponent policies, augment the population, and compute meta-strategies in the resulting empirical game. The PSRO paper emphasizes mitigating overfitting to a single opponent via meta-game training.

% Population Based Training (PBT) is a practical hyperparameter-and-weight evolution method that can be layered on top of MARL training to maintain diversity and adapt learning rates/entropy bonuses online.

% \subsubsection{Sequence-based policies for theorem-proving-style VAMP observations}

% VAMP observations are variable-length structured objects (libraries, dependency graphs, market ledgers). The draft proposes a causal transformer actor over tokenized trajectories (a decision-transformer-like architecture).

% Decision Transformer reframes RL as conditional sequence modeling on (return, state, action) token sequences, enabling transformer-based policies over long contexts.

% In our setting, we use the transformer primarily as a \textbf{policy network class} handling long action--observation histories and variable-length candidate sets (e.g., pointer heads over entities), while PPO/MAPPO supplies the on-policy improvement loop. (This hybrid ``transformer actor + PPO'' design is increasingly common in practice; formal theory is limited. \textbf{TODO:} cite a primary study of PPO with transformer policies in partially observable settings.)

% \subsection{Convergence Theory and Provable Guarantees}

% We separate guarantees that are genuinely theorem-level from those that are heuristic.

% \subsubsection{No-regret implies CCE in repeated normal-form (meta-)games}

% The foundational guarantee, used in your Chapter 7 draft, is that no-external-regret dynamics converge (in time-average) to coarse correlated equilibrium in repeated play.

% \paragraph{Definition (CCE).}
% For a finite normal-form game with action profiles \(a\in\mathcal A=\prod_\alpha \mathcal A^\alpha\), a distribution \(\mu\in\Delta(\mathcal A)\) is a coarse correlated equilibrium if for every agent \(\alpha\) and every fixed deviation action \(\tilde a^\alpha\in\mathcal A^\alpha\),
% \[
% \mathbb E_{a\sim\mu}[u^\alpha(a)] \;\ge\; \mathbb E_{a\sim\mu}[u^\alpha(\tilde a^\alpha,a^{-\alpha})].
% \]
% (For correlated equilibrium proper, one uses \textbf{internal regret}; see Foster--Vohra / Hart--Mas-Colell.)

% \paragraph{Theorem (external regret \(\Rightarrow\) approximate CCE).}
% Consider repeated play for \(T\) rounds in a fixed normal-form game. Let \(\hat\mu_T\) be the empirical distribution over joint actions. If every player \(\alpha\) has external regret \(\mathrm{Reg}^\alpha(T)\le \varepsilon T\), then \(\hat\mu_T\) is an \(\varepsilon\)-CCE:
% \[
% \forall \alpha,\forall \tilde a^\alpha:\quad
% \mathbb E_{a\sim\hat\mu_T}[u^\alpha(a)] \ge \mathbb E_{a\sim\hat\mu_T}[u^\alpha(\tilde a^\alpha,a^{-\alpha})] - \varepsilon.
% \]

% \emph{Proof.} Fix \(\alpha\) and \(\tilde a^\alpha\). External regret bound implies
% \[
% \sum_{t=1}^T u^\alpha(a_t)\ \ge\ \sum_{t=1}^T u^\alpha(\tilde a^\alpha,a_t^{-\alpha}) - \varepsilon T.
% \]
% Divide by \(T\) and recognize empirical expectations under \(\hat\mu_T\). \(\square\)

% This is the precise mathematical sense in which ``no-regret learning converges to CCE.'' It is the load-bearing convergence statement we can safely transport into our VAMP-MARL framework, \emph{but only at the level of a normal-form meta-game}.

% \subsubsection{From VAMP to a policy meta-game}

% Define a finite population \(\Pi^\alpha=\{\pi^{\alpha,1},\dots,\pi^{\alpha,m_\alpha}\}\). The induced \textbf{empirical meta-game} has ``actions'' = choosing an index in \([m_\alpha]\) for each agent; payoff is expected episodic return of the corresponding joint policy in the underlying VAMP. This is standard in empirical game-theoretic analysis and PSRO.

% If our training loop ensures that each agent's meta-strategy updates have low external regret w.r.t.\ \(\Pi^\alpha\), then the empirical distribution over played policies converges to an approximate CCE of the meta-game by the theorem above. This is exactly the ``weak convergence'' guarantee described in the draft: it is a guarantee about the population game, not the full POSG/VAMP.

% \subsubsection{PSRO guarantees in two-player zero-sum}

% In two-player zero-sum settings, PSRO with exact best-response oracles and an exact meta-solver converges (in the limit) to a Nash equilibrium of the underlying game; the PSRO paper frames this as a unifying view of double-oracle/fictitious play style methods. In VAMP, this gives a principled \emph{sanity-check regime} (fully observable, two-player, zero-sum markets) where equilibrium approximation is theoretically grounded.

% \textbf{TODO:} state the precise convergence theorem you want to cite (e.g., double-oracle convergence in finite zero-sum games) and map it to your PSRO instantiation.

% \subsubsection{Impossibility / lower bounds inherited from PPAD and NEXP}

% Our equilibria chapter already proves that general equilibrium computation is intractable: PPAD-hard even for one-shot subcases, and NEXP-hard for cooperative partial observability. Consequently, any \emph{general-purpose} MARL method with a polynomial-time worst-case guarantee to output (even approximate) Nash for arbitrary VAMPs would imply breakthroughs in PPAD-hard problems; similarly, any method that guarantees optimal cooperative performance for arbitrary cooperative VAMPs would imply breakthroughs against NEXP-hard Dec-POMDP planning.

% We therefore frame MARL guarantees as: convergence to approximate equilibrium notions \textbf{within a restricted class} (policy parameterizations + population), under \textbf{stability assumptions} (mixing, bounded gradients, sufficient exploration), none of which hold uniformly in worst-case VAMPs.

% \subsection{Practical Algorithm Design for VAMPs}

% This section instantiates the draft's design choices in a principled MARL lens.

% \subsubsection{Credit assignment across task and market layers}

% VAMP rewards combine: (i) task utility (resolved-node value, unlock effects) and (ii) market P\&L plus collateral costs. This creates long-horizon credit assignment and counterfactual dependence (an agent's publication changes others' feasible sets and prices).

% Practical choices:
% \begin{itemize}
% \item \textbf{Centralized critic features:} include global graph statistics (depth, number of unresolved prerequisites), market summaries (best bids/asks, AMM price vector), and public-library deltas.
% \item \textbf{Reward decomposition (optional):} shape rewards into immediate components (e.g., bounty capture, duplication penalty, unlock bonus). \textbf{TODO:} cite potential-based shaping conditions ensuring policy invariance in MDPs; adaptation to POSGs is nontrivial.
% \item \textbf{Counterfactual baselines (optional):} COMA-style counterfactual advantage can reduce variance for discrete multiagent actions. \textbf{TODO:} cite COMA primary paper if you adopt it explicitly.
% \end{itemize}

% \subsubsection{Exploration under endogenous task creation}

% Conjecture actions expand the task graph; naive exploration can explode the state. We use curriculum learning (small DAGs \(\to\) deeper DAGs, few agents \(\to\) many agents, simple markets \(\to\) rich mechanisms), as in the draft.

% Exploration techniques suited to VAMP:
% \begin{itemize}
% \item \textbf{Action masking:} forbid inadmissible trades (collateral) and logically impossible task actions.
% \item \textbf{Structured entropy:} separate entropy bonuses for action-type, target-selection, and parameter heads to keep exploration in each factorized component.
% \item \textbf{Opponent sampling:} population-based opponent mixtures to avoid brittle overfitting (PSRO motivation).
% \end{itemize}

% \subsubsection{Handling market actions and collateral constraints}

% Market actions are constrained by admissibility; we treat admissibility as a hard constraint via masking. For AMMs, portfolio changes are deterministic given trade parameters; we can incorporate the resulting P\&L into immediate reward shaping to reduce variance.

% \paragraph{Unspecified detail.}
% Exact collateral rule (fully collateralized vs.\ margin). Reasonable options: (i) fully collateralized (simplifies masking), (ii) bounded leverage with forced liquidation (adds nonstationary constraints). This choice materially changes exploration and critic stability.

% \subsection{Training Pipeline, Baselines, and Protocol}

% We adopt the staged pipeline already drafted.

% \begin{enumerate}
% \item \textbf{Offline initialization (behavior cloning):} supervised learning on scripted baseline trajectories to reduce cold-start exploration.
% \item \textbf{Online CTDE fine-tuning (MAPPO/PPO):} rollout \(\to\) GAE advantage estimation \(\to\) clipped PPO updates \(\to\) critic updates.
% \item \textbf{Population-based training:} maintain diverse roster (current, checkpoints, scripted) and sample opponents; optionally evolve hyperparameters with PBT.
% \item \textbf{Curriculum:} scale instance difficulty along graph size, depth, number of agents, market richness, and information asymmetry.
% \end{enumerate}

% \paragraph{Baselines (for evaluation and ablation).}
% Centralized oracle (upper bound), assignment/LP snapshot heuristics, dependency-aware list/priority rules, price-taking heuristics, and ``no-market'' VAP agents.

% \paragraph{Metrics.}
% Resolved utility, duplication avoidance (wasted compute), market efficiency (spread/liquidity utilization), regret/exploitability in restricted meta-games, and generalization across market mechanisms.

% \subsection{Illustrative Diagrams and Tables}

% \subsubsection{Training loop diagram}

% \begin{figure}[t]
% \centering
% \fbox{\begin{minipage}{0.9\textwidth}
% \centering
% \textbf{VAMP simulator}\\
% \(\downarrow\) rollouts \((w_t, o_t^\alpha, a_t, r_t)\)\\
% \textbf{Trajectory buffer}\\
% \(\swarrow\) \hspace{2cm} \(\searrow\)\\
% \textbf{Centralized critic \(V_\psi(w_t)\)} \hspace{1cm} \textbf{Decentralized actors \(\pi_\theta^\alpha(h_t^\alpha)\)}\\
% \(\downarrow\) advantages \(A_t\) \hspace{1cm} \(\downarrow\) policy snapshots\\
% \textbf{PPO/MAPPO updates} \hspace{1cm} \textbf{Policy population / league}\\
% \hspace{5.2cm} \(\downarrow\) opponent sampling\\
% \hspace{5.2cm} \textbf{VAMP simulator}\\
% \textbf{Scripted baselines} \(\rightarrow\) \textbf{Policy population / league}
% \end{minipage}}
% \caption{Training loop for VAMP-MARL.}
% \label{fig:vamp-marl-training-loop}
% \end{figure}

% \subsubsection{Algorithm comparison table}

% \begin{table}[t]
% \centering
% \small
% \begin{tabular}{lcccc}
% \hline
% Method & CTDE? & Sample eff. & Stability & VAMP suitability \\
% \hline
% MAPPO (PPO+central critic) & yes & medium & high & strong default \\
% MADDPG (off-policy CTDE) & yes & high & medium/low & good for continuous \\
% PSRO / league training & n/a & medium & high & strong for robustness \\
% Decision Transformer (offline) & no/optional & high (offline) & medium & strong for long contexts \\
% PBT (wrapper) & n/a & n/a & improves & good for hyperparams/diversity \\
% \hline
% \end{tabular}
% \caption{Qualitative MARL options for VAMP training.}
% \end{table}

% Sources for these algorithm families: PPO, MAPPO-style PPO in cooperative games, CTDE actor--critic (MADDPG), PSRO, Decision Transformer, and PBT.

% \subsection{Prioritized Bibliography and Open Problems}

% Prioritized sources to cite in this chapter section: \textbf{Reinforcement Learning: An Introduction} (Sutton--Barto 2018), \textbf{John Schulman} on PPO, MAPPO empirical study, \textbf{Ryan Lowe} on CTDE actor--critic, \textbf{Marc Lanctot} on PSRO, \textbf{Lili Chen} on Decision Transformer, \textbf{Max Jaderberg} on PBT, and no-regret foundations: \textbf{Sergiu Hart} / \textbf{Andreu Mas-Colell}, plus online learning text (Cesa-Bianchi--Lugosi).

% Open problems (tight to later implementation/experiments): (i) formalize regret in continuous policy-parameter spaces with function approximation for VAMP meta-games; (ii) quantify how collateral constraints change exploration and stability; (iii) characterize when CTDE critics actually reduce nonstationarity in market-mediated POSGs; (iv) derive mechanism-specific learning dynamics (e.g., AMM price response + agent gradients). These connect directly to the implementation and experiments chapters described in the thesis outline.
